%%%%%%%%%%%%%%%%%%%%%%%%%%%%%%%%%%%%%%%%%%%%%%%%%%%%%%%
% BivariateAggregate Cheat Sheet
%
% Created by Stephen J Mildenhall
% (c) 2026
%
% New at aggregate v1.0.0a169 (REFACTOR branch):
%   * content regenerated from introspect(build('bv ...')) and from
%     introspect(bv.bivariate) (the BivariateDistribution view)
%   * four declaration modes: copula, discrete (dbvsev), view-pair
%     (netceded / grossceded / grossnet), and clash
%   * massive (disk-backed) update via update(store_dir=...)
%
%%%%%%%%%%%%%%%%%%%%%%%%%%%%%%%%%%%%%%%%%%%%%%%%%%%%%%%

\input{cheat_sheet_macros.tex}
\title{BivariateAggregate Cheat Sheet}

% color scheme defeined here - just change the suffixes
\colorlet{highlightcolor}{highlightcolorf}
\colorlet{washedcolor}{washedcolorf}
\colorlet{textcolor}{textf}


\begin{document}

{\huge{\textbf{\texttt{BivariateAggregate} Class Cheat Sheet}}}

\raggedright
\texttt{\m BivariateAggregate(name, units=None, copula=None, note='', hints='', tags=(), doc='', mode='copula', nc\_agg, nc\_kwargs, nc\_views,} \\
\texttt{\phantom{\m}clash, dbv\_xs, dbv\_ys, dbv\_S, exp\_en, exp\_el, exp\_premium, exp\_lr, freq\_name='poisson', freq\_a, freq\_b, freq\_zm, freq\_p0)}

A \texttt{BivariateAggregate} is the joint law $(A_0, A_1)$ of two coupled aggregates. A \textbf{shared} outer frequency drives both; within each event the two per-claim severities are coupled, and the joint is accumulated by a 2D FFT. Either marginal reproduces the standalone \texttt{Aggregate}. Usually built via \texttt{build('bivariate ...')} (alias \texttt{bv}); see the DecL card, page 3.
The following tables show all \texttt{\m methods}, and fields or properties (used interchangeably). Comments elucidate the meaning of more obscure entries.


\begin{multicols*}{3}


%------------1. SPECIFICATION & CREATION ---------------
\begin{tikzpicture}
\node [mybox] (box){%
    \begin{minipage}{0.3\textwidth}\raggedright

\texttt{name},
\texttt{units} (the two component \texttt{Aggregate}s),
\texttt{unit\_names},
\texttt{copula} (the fitted \texttt{Copula}),
\texttt{frequency}, \texttt{freq\_name} (the \emph{shared} outer count),
\texttt{en} (expected event count),
\texttt{clash} (the solved $(n_a, n_b, n_c, n_0, p_a, p_b)$ when declared with \texttt{clash}),
\texttt{program}, \texttt{pprogram},
\texttt{mode}${}^{[1]}$

\smallskip
{\it Four declaration modes:}
\texttt{copula} (two \texttt{agg} bodies $+$ a \texttt{copula} clause),
\texttt{discrete} (\texttt{dbvsev}, the joint per-claim matrix given directly),
view-pair (\texttt{netceded} / \texttt{grossceded} / \texttt{grossnet} over one reinsured \texttt{agg}),
\texttt{clash} (counts $n_a$, $n_b$, $n_c$ solved into a shared count and two Bernoulli triggers).


    \end{minipage}
};
\node[fancytitle] at (box.north west) {1. Specification \& creation};
\end{tikzpicture}


%------------2. UPDATE ---------------
\begin{tikzpicture}
\node [mybox] (box){%
    \begin{minipage}{0.3\textwidth}\raggedright

\texttt{\m update(log2, bs, padding, store\_dir, row\_chunk, col\_chunk, keep\_transform)},
\texttt{\m update\_work},
\texttt{bs} (per-axis bucket sizes, a \emph{list}),
\texttt{axis\_xs} (per-axis output grids),
\texttt{padding}

\smallskip
{\it Scalar \texttt{log2} is the \rm total\it{} 2-D cell budget, split between the axes by measured support; a 2-tuple pins the per-axis \texttt{log2}. \texttt{store\_dir} switches to the \rm massive\it{} disk-backed path (extra \texttt{aggregate[massive]}), where \texttt{padding} defaults to 0.}


    \end{minipage}
};
\node[fancytitle] at (box.north west) {2. Update};
\end{tikzpicture}


%------------ 3. MOMENTS ---------------
\begin{tikzpicture}
\node [mybox] (box){%
    \begin{minipage}{0.3\textwidth}\raggedright

\texttt{\m moments} (mixed raw moments $\mathsf{E}[A_0^i A_1^j]$),
\texttt{corr} (Pearson correlation of the two aggregates),
\texttt{dependency\_df} (\texttt{cov} / \texttt{corr} / \texttt{tau} at both levels: per-claim \texttt{Sev} and \texttt{Agg}),
\texttt{stats\_df} (per-component marginal moments, theory vs.\ empirical)


    \end{minipage}
};
\node[fancytitle] at (box.north west) {3. Moments};
\end{tikzpicture}


\columnbreak


%------------ 4. STATISTICAL FUNCTIONS ---------------
\begin{tikzpicture}
\node [mybox] (box){%
    \begin{minipage}{0.3\textwidth}\raggedright

{\it Via the \rm\texttt{bivariate}\it{} view}${}^{[2]}$:
\texttt{\m pushforward} (law of $z = f(A_0, A_1)$ on a fresh 1-D grid --- the workhorse: \emph{any} scalar function of the pair),
\texttt{\m transformed\_moments} (exact raw and central moments of $f(A_0, A_1)$, no regridding),
\texttt{\m moments}, \texttt{\m corr}, \texttt{\m marginals}

\smallskip
{\it Univariate risk measures come from the pushforward or from either marginal \texttt{Aggregate}; a 2-D law has no single quantile.}


    \end{minipage}
};
\node[fancytitle] at (box.north west) {4. Statistical functions};
\end{tikzpicture}


%------------ 5. VALIDATION ---------------
\begin{tikzpicture}
\node [mybox] (box){%
    \begin{minipage}{0.3\textwidth}\raggedright

\texttt{summary\_df} (two-unit \texttt{Portfolio}-shape summary: theory vs.\ realised),
\texttt{validation\_explanation},
\texttt{deficit} (mass lost off the grid corner)${}^{[3]}$,
\texttt{bs\_window\_df}, \texttt{bs\_description}, \texttt{bs\_explanation} (per axis),
\texttt{tail\_df}, \texttt{tail\_description}, \texttt{tail\_explanation}


    \end{minipage}
};
\node[fancytitle] at (box.north west) {5. Validation};
\end{tikzpicture}


%------------ 6. OUTPUT ---------------
\begin{tikzpicture}
\node [mybox] (box){%
    \begin{minipage}{0.3\textwidth}\raggedright

\texttt{density} (the joint pmf, shape $(n_0, n_1)$),
\texttt{density\_df} (the same as a \texttt{DataFrame}),
\texttt{marginals} (the two 1-D densities),
\texttt{bivariate} (the \texttt{BivariateDistribution} view)${}^{[2]}$,
\texttt{info},
\texttt{summary\_df}, \texttt{stats\_df}, \texttt{dependency\_df},
\texttt{axis\_xs}, \texttt{figure}


    \end{minipage}
};
\node[fancytitle] at (box.north west) {6. Output dataframes};
\end{tikzpicture}


%------------ 7. REINSURANCE ---------------
\begin{tikzpicture}
\node [mybox] (box){%
    \begin{minipage}{0.3\textwidth}\raggedright

{\it The whole point of the view-pair modes.} \texttt{netceded <AGG>} (and \texttt{grossceded}, \texttt{grossnet}) build the joint occurrence law of two of \{gross, ceded, net\} for one reinsured aggregate --- exact dependence, comonotone within the layer. The object-API route is \texttt{\m Aggregate.occ\_bivariate(views)}.

\smallskip Per-component reinsurance is declared on each inner \texttt{agg} line as usual.


    \end{minipage}
};
\node[fancytitle] at (box.north west) {7. Reinsurance};
\end{tikzpicture}


\columnbreak


%------------ 8. VISUALIZTION ---------------
\begin{tikzpicture}
\node [mybox] (box){%
    \begin{minipage}{0.3\textwidth}\raggedright

\texttt{\m plot} (two-panel contour: per-claim severity, then aggregate),
\texttt{\m bivariate.contour} (one filled contour panel),
\texttt{figure}


    \end{minipage}
};
\node[fancytitle] at (box.north west) {8. Visualization};
\end{tikzpicture}


%------------ 9. RISK ---------------
\begin{tikzpicture}
\node [mybox] (box){%
    \begin{minipage}{0.3\textwidth}\raggedright

{\it None directly.} Push the joint through the payoff you want to price --- \texttt{bivariate.pushforward(f)} returns an ordinary 1-D distribution --- then price that, or wrap it in a \texttt{PnL} (a \texttt{BivariateDistribution} is an accepted P\&L \texttt{source}).


    \end{minipage}
};
\node[fancytitle] at (box.north west) {9. Risk and pricing};
\end{tikzpicture}


%------------  10. APPROXIMATIONS ---------------
\begin{tikzpicture}
\node [mybox] (box){%
    \begin{minipage}{0.3\textwidth}\raggedright

{\it None.} The 2-D FFT is exact on the grid; the only approximation levers are the axis budgets (\texttt{log2}, \texttt{bs}) and \texttt{padding}, in Update.

    \end{minipage}
};
\node[fancytitle] at (box.north west) {10. Approximations};
\end{tikzpicture}


%------------ 11. META ---------------
\begin{tikzpicture}
\node [mybox] (box){%
    \begin{minipage}{0.3\textwidth}\raggedright

\texttt{\m help},
\texttt{\m format\_program}, \texttt{pprogram\_html},
\texttt{note}, \texttt{tags}, \texttt{doc}, \texttt{hints} (the DecL trailer)

    \end{minipage}
};
\node[fancytitle] at (box.north west) {11. Meta};
\end{tikzpicture}

\bigskip \raggedright
{\bf Notes:}

[1]: \texttt{mode} is \texttt{'copula'} or \texttt{'discrete'}; the view-pair and \texttt{clash} statements set the remaining constructor arguments (\texttt{nc\_agg} / \texttt{nc\_views}, \texttt{clash}).

[2]: \texttt{bivariate} is a \texttt{BivariateDistribution}: \texttt{density}, \texttt{axis0} / \texttt{axis1} (positionally \texttt{ceded} / \texttt{net}), \texttt{axis\_names}, \texttt{bs\_ceded} / \texttt{bs\_net}, \texttt{meta}, and \texttt{\m marginals}, \texttt{\m moments}, \texttt{\m corr}, \texttt{\m contour}, \texttt{\m pushforward}, \texttt{\m transformed\_moments}. Address axes \emph{by role} (\texttt{axis0} / \texttt{axis1} / \texttt{axis\_names}), never by the positional names.

[3]: With \texttt{padding=0} the FFT wraps rather than losing mass, so \texttt{deficit} can read clean while the answer is wrong: compare means, not \texttt{deficit}.


% \vskip .25truein
% FOOTER
\makefooter



\end{multicols*}

\end{document}
